\chapter{The Legacy Compiler Compatibility Protocol}\label{ch07:worker-protocol}

\section{The profile and its refusal}\label{ch07:profile}

Every mechanism described so far in this monograph---structural semantic references, the research boundary object, the canonical evidence encoding, the byte-backed registries---concerns artifacts at rest. This chapter concerns the one place in CAS Core 8 where something actually \emph{runs}: the legacy compiler compatibility profile, under which the kernel executes the frozen CAS v7.7 Python notation compiler against caller-supplied source text and returns a hash-bound record of what happened.

The profile is named \texttt{legacy\_only}, and the name is a commitment, not a description. The profile binds the exact bytes of the active Python CAS-notation compiler---the single compiler artifact listed in the Gate~4 record, identified by a raw digest and a newline-normalized content digest---and it executes those bytes, and only those bytes, for comparison purposes. What it refuses to do is equally definite: it produces \emph{no} CAS v8 semantic translation. The compilation record carries a field \texttt{v8\_semantic\_translation} whose only admissible value is the null value; a constructor that receives anything else raises \texttt{E-COMPAT-V8-TRANSLATION} and the record never exists. There is no configuration, flag, or subclass through which a legacy compilation acquires v8 meaning. If a future v8 semantics wants to say something about a legacy program, it must do so through a separately reviewed translation gate that does not yet exist, and the compatibility record is deliberately shaped so that its output cannot be mistaken for that gate's output.

This refusal is the first thing to understand about the protocol, because it explains the protocol's otherwise puzzling austerity. A compatibility execution is not trying to be useful in the ordinary sense---it is not an interpreter service, a build system, or a language server. It is an \emph{evidence procedure}: a way of establishing, to the standard the kernel's canonical encoding can support, the statement ``these exact source bytes, compiled by these exact compiler bytes, on this exactly described interpreter, produced this exact artifact.'' Every step below exists because some component of that statement can be silently falsified in an ordinary execution environment, and the protocol's job is to close each falsification channel or to record, honestly, that it remains open.

\begin{definition}[The frozen compiler identity]\label{ch07:def:compiler-identity}
The admitted compiler is identified by three constants: a raw digest over the on-disk bytes as they exist in the working tree (historically carrying CRLF line endings, digest beginning \texttt{065DC438}), a canonical digest over the LF-normalized content (digest beginning \texttt{8169188A}), and a source-authority identifier \texttt{SRC-ACTIVE-COMPILER-CURRENT}. The identity policy accepts exactly two raw variants---the CRLF form and the LF form---and requires that the LF-normalized content digest equal the single frozen canonical value in every case.
\end{definition}

The two-variant acceptance is not laxity; it is the byte-identity discipline of the canonical encoding chapter applied to a real repository. A checkout on Windows and a checkout on a POSIX host may legitimately hold different raw bytes for the same compiler because of newline rewriting, and the program has been burned before by conflating the two identities. The policy therefore admits either raw form but anchors both to one LF content identity, and---critically---the receipt records \emph{both} digests separately, so a reviewer can always tell which physical variant actually executed. Neither digest substitutes for the other, and no other raw variant is admissible: a compiler file whose LF normalization matches but whose raw digest is neither frozen variant is rejected, because an unexplained third physical form is itself evidence of something the profile does not understand.

\begin{remark}
Throughout this chapter, error identifiers such as \texttt{E-COMPAT-SOURCE-HASH} name the typed failure codes raised by the reference implementation. They are quoted because the taxonomy is load-bearing: Chapter~\ref{ch08:concurrency} argues that \emph{which} failure occurred must survive into evidence, and that requires failures to be enumerable rather than prose.
\end{remark}

The remainder of the chapter develops the eight-step deterministic-closure protocol, one section per step. The protocol statement itself is short; the argument for each step is where the content lives.

\begin{protocol}[Deterministic closure, restated]\label{ch07:prot:closure}
A compatibility execution: (1) verifies the requested source hash and reads and hashes the compiler bytes before execution; (2) forces the Python-only engine profile under an exact frozen policy; (3) rejects \texttt{use} imports, because no import-DAG manifest is admitted; (4) executes the verified compiler bytes in a fresh \texttt{python -I -S -B} worker with a sanitized environment; (5) samples and binds the worker runtime before compiler code executes, without WMI or subprocess probing; (6) serializes worker launches through a profile-bound advisory lock with a 90-second acquisition bound, 10-millisecond poll, and 10-millisecond post-release yield, never retrying a failed worker and never claiming fairness; (7) records the full receipt---raw and LF compiler identities, interpreter and worker bytes, import search paths, loaded-module inventory, exact syntax, semantic, and engine profiles, result digest, diagnostics, and research boundary; and (8) requires fresh replay against the same source and compiler identity before any structural decision can pass.
\end{protocol}

\section{Step 1: source-hash verification and pre-execution compiler-byte hashing}\label{ch07:step1}

The first step establishes \emph{what} is about to run, before anything runs.

The caller does not hand the protocol a file path and a promise. Source text arrives inside a typed envelope carrying the source string, a declared version, a declared syntax profile, and a SHA-256 digest of the source's UTF-8 encoding. The protocol's first act is to recompute that digest from the retained source bytes and compare it to the declared value; a mismatch raises \texttt{E-COMPAT-SOURCE-HASH} and nothing executes. The declared version must be one of the two accepted version strings and the declared syntax profile must be exactly \texttt{cas-v7.7}; anything else fails closed with \texttt{E-COMPAT-VERSION} or \texttt{E-COMPAT-PROFILE}.

Symmetrically, the compiler is read from disk \emph{once}, as bytes, and both of its identities---raw and LF-normalized---are computed from those retained bytes before the worker is launched. The retained bytes, not the path, are what the worker later receives: the compiler source is transmitted to the child process over standard input, base64-encoded inside the worker request, together with its raw digest. The worker independently recomputes the digest of what it received and refuses to execute on mismatch. The path from which the bytes were read plays no further role in the evidence.

\begin{threat}[Time-of-check to time-of-use substitution]\label{ch07:threat:toctou}
An adversary---or, far more commonly, an innocent concurrent process such as an editor, a formatter, or a version-control operation---modifies the compiler file between the moment its hash is computed and the moment it is executed. A protocol that hashes a path and then executes the path certifies the wrong bytes.
\end{threat}

Step 1 blocks this threat structurally rather than by racing faster. Because the bytes that are hashed are the same in-memory bytes that are serialized into the worker request, there is no window in which the file system can substitute content: the file could be rewritten a nanosecond after the read and the evidence would still be exactly correct, because the evidence describes the retained bytes, not the file. The compilation record then retains those compiler bytes for its whole lifetime, and its constructor re-derives the identity from the retained bytes and compares it against the recorded identity---so even a record deserialized or reconstructed later cannot claim an identity its own retained bytes contradict (\texttt{E-COMPAT-COMPILER-IDENTITY}).

The same pattern applies to the source. Note the doubled check: the envelope verifies its own digest at construction, and the compilation constructor re-verifies both that the envelope's source still hashes to the envelope's declared digest and that the compilation's top-level \texttt{source\_sha256} equals the envelope's. Redundant verification of this kind is a deliberate stylistic commitment of the kernel. Each object recomputes, from primitive retained data, every identity it asserts, rather than trusting that some earlier layer already did. The cost is a few extra hash computations; the benefit is that no single bypassed check silently propagates a false identity into evidence.

Why hash \emph{before} execution rather than after? Because a post-execution hash certifies at most what the file looked like after the run, and an execution can modify its own inputs---or an adversary can restore a tampered file after the run completes. Pre-execution hashing over retained bytes is the only ordering under which the digest provably describes what executed. This is the reproducible-builds discipline \cite{reprobuilds} applied at the granularity of a single tool invocation: the artifact's identity is fixed at the moment of capture, and everything downstream refers to the capture, not to the mutable world.

\begin{ruleenv}[No path authority]\label{ch07:rule:no-path-authority}
A file path is a lookup hint, never an identity. Every identity claim in the compatibility profile is over retained bytes, and every consumer of the claim can recompute it from data the record itself carries.
\end{ruleenv}

\section{Step 2: the forced Python-only engine profile and the frozen policy}\label{ch07:step2}

The second step establishes \emph{what kind of thing} is allowed to do the running.

The kernel admits exactly one engine profile in v0.1: \texttt{cas-v7.7/python-only}. The constant is not a default that a caller may override; it is checked at three independent points. First, the frozen policy object---compiler source authority, accepted versions, accepted syntax profiles, both raw digest variants, and the canonical LF digest---is reconstructed at every compilation and verified field-by-field against the module constants; if the policy object and the constants have drifted apart in any component, the compilation fails with \texttt{E-COMPAT-POLICY} before any worker launches. Second, the worker's returned artifact must itself declare \texttt{engine} equal to \texttt{python}; a legacy compiler that dispatched to any other execution engine, however it might have been induced to, fails the run with \texttt{E-COMPAT-ENGINE}. Third, the compilation record's constructor refuses any engine profile string other than the frozen constant, so the restriction survives into every downstream copy of the evidence.

What failure does this triple check block? Three distinct ones.

\emph{Configuration drift.} A policy expressed once, at a distance from its enforcement point, tends to decay: constants get edited for a test and not reverted, a refactor renames a profile string, a merge resolves a conflict wrongly. By reconstructing and re-verifying the policy at each compilation, the protocol converts silent drift into a loud, typed failure at the next execution. The policy is not configuration in the operational sense at all---there is nothing to configure---it is a self-checking restatement of the profile's identity.

\emph{Engine substitution.} The CAS program contains, and intends to contain, more than one execution engine: the monorepo's architecture designates a Rust engine as the eventual production scoring core, and cross-implementation conformance between Python and Rust is explicitly \emph{not established}---it is an open obligation, awaiting a differential-testing campaign. Until that campaign exists, a receipt that said only ``the compiler ran'' would be ambiguous across engines whose agreement no one has demonstrated. Forcing the Python-only profile makes the receipt's scope exact: this evidence speaks about the Python execution path and nothing else. A future native-engine profile would be admitted, if at all, as a \emph{separately hash-bound} profile with its own identity constants---never as a widening of this one.

\emph{Semantic scope creep.} The semantic profile string recorded in every receipt is \texttt{CAS v7.7 operational compatibility only}. The phrase is chosen to be un-misreadable: the profile asserts operational compatibility---the frozen compiler, run under the frozen policy, produced this artifact---and nothing about the correctness, meaning, or v8 interpretation of the compiled notation. A receipt is not a semantics. Freezing the semantic profile string, and rejecting receipts that carry any other, prevents the gradual rewording by which scope claims usually grow.

\begin{invariant}[One profile, three checks]\label{ch07:inv:one-profile}
In v0.1 the engine profile is a single frozen constant verified independently at policy reconstruction, at worker-artifact inspection, and at record construction. No admitted code path executes legacy source under any other engine profile, and no admitted record asserts one.
\end{invariant}

\section{Step 3: rejection of \texttt{use} imports}\label{ch07:step3}

The third step bounds \emph{how much} is allowed to run: exactly one source unit.

The CAS notation language includes a \texttt{use} construct by which one program imports another. The compatibility profile rejects it entirely. After the worker returns, the protocol walks the returned abstract syntax tree and fails the compilation with \texttt{E-COMPAT-UNBOUND-IMPORT} if any node of kind \texttt{UseStmt}, \texttt{UseStatement}, or \texttt{ModuleImport} appears anywhere in it. Multi-file CAS programs are, in v0.1, simply outside the profile.

The reason is not that imports are dangerous in themselves; it is that the profile has no honest way to \emph{identify} what an import would bring in. Every other input to a compatibility execution is content-addressed: the source is hash-verified, the compiler is hash-verified, the interpreter and worker and loaded Python modules are inventoried and digested. An import statement, by contrast, names a module by a resolvable name, and resolution is a function of search paths, file-system state, and timing. To admit imports honestly, the profile would need an \emph{import-DAG manifest}: a caller-supplied, hash-bound closure enumerating every transitively imported source unit with its digest, verified before execution exactly as the root source is, so that the identity ``this program'' expands to ``this exact finite set of byte strings.'' No such manifest format has been designed, reviewed, or admitted. Until one is, an execution containing an import would be an execution whose full input is unidentified---and an evidence record over an unidentified input is precisely the kind of artifact this kernel exists to make unrepresentable.

\begin{threat}[Ambient-resolution smuggling]\label{ch07:threat:ambient-import}
A source program imports a module whose content differs between the compilation that produced a receipt and a later replay, or differs between two hosts, or is chosen adversarially via search-path manipulation. The receipt appears to bind ``the program'' while an unbounded portion of the program's meaning rides in un-hashed ambient state.
\end{threat}

The rejection blocks this threat at the grammar level, which is the strongest available place: no analysis of resolution behavior is needed, because the construct that would invoke resolution is not admitted. Notice also \emph{where} the check runs---on the returned AST, after compilation---rather than by pre-scanning the source text. Textual pre-scanning would have to reimplement enough of the lexer to distinguish a real \texttt{use} statement from the word appearing inside a string or comment, and any divergence between the scanner and the actual compiler grammar becomes a bypass. Inspecting the compiler's own parse output means the check sees imports exactly where the compiler sees them.

The AST walk is complemented by a second, independent emptiness check: the worker reports an \texttt{external\_cas\_imports} list in its isolation record, and both the protocol and the execution-profile constructor require that list to be empty (\texttt{E-COMPAT-UNBOUND-IMPORT} again). Two mechanisms, one from the artifact side and one from the runtime side, must agree that nothing external was touched.

\begin{boundary}[Out of profile, not impossible]\label{ch07:bd:multifile}
The rejection of multi-file programs is a scope statement, not a capability judgment. Admitting an import-DAG manifest is listed among the kernel's open obligations; nothing in this chapter should be read as claiming that multi-file compatibility evidence cannot be built---only that it has not been, and that the profile fails closed rather than approximating it.
\end{boundary}

\section{Step 4: the fresh \texttt{python -I -S -B} worker and the sanitized environment}\label{ch07:step4}

The fourth step controls \emph{where} the run happens: a freshly created Python process whose interpreter-level state is as close to a function of its declared inputs as CPython permits.

The worker is launched as \texttt{python -I -S -B -u} on the resolved path of the currently selected interpreter, with the packaged worker script as its program, the request delivered over standard input, a freshly created temporary directory as its working directory, a 30-second wall-clock timeout, and an environment constructed from scratch rather than inherited. Each flag closes a specific channel.

\emph{Isolated mode, \texttt{-I}.} This is the broadest of the three flags and subsumes several older ones. It removes the script's own directory (and the current working directory) from the module search path; it ignores all \texttt{PYTHON*} environment variables that would otherwise steer the interpreter (\texttt{PYTHONPATH}, \texttt{PYTHONSTARTUP}, and the rest); and it disables the per-user site-packages directory. The channels it closes are exactly the classic interpreter-hijack channels: a \texttt{PYTHONPATH} entry that shadows a standard-library module, a startup file that runs arbitrary code before the worker's first line, a user-site package that monkey-patches \texttt{json} or \texttt{hashlib}. Under \texttt{-I}, the interpreter's import behavior is a function of the installation, not of the invoking user's ambient configuration.

\emph{No site initialization, \texttt{-S}.} Even without user-site directories, CPython ordinarily imports the \texttt{site} module at startup, which executes \texttt{sitecustomize} and any \texttt{.pth} files found in site directories---and \texttt{.pth} files may contain executable lines. \texttt{-S} skips site initialization entirely. This matters for evidence because \texttt{.pth}-injected code is close to invisible: it runs before any user code, appears in no import statement, and varies per installation. A worker that ran site initialization could have its module namespace altered by content the receipt never inventoried.

\emph{No bytecode writing, \texttt{-B}.} The worker must not write \texttt{.pyc} files. This is partly hygiene---the worker should not leave residue keyed to its inputs on the shared host---but it is also an evidence measure: cached bytecode is a second physical representation of source whose staleness semantics are timestamp-based, and a future run that silently loaded a stale cache would execute bytes that no receipt hashed. Refusing to write the cache in the worker removes the worker as a source of such artifacts. (The unbuffered flag \texttt{-u} is operational rather than evidential: it ensures the single JSON response is not stranded in a stdio buffer when the process exits.)

The environment is rebuilt, not filtered. The construction starts from an empty dictionary and copies in an enumerated whitelist: on all platforms \texttt{SystemRoot}, \texttt{WINDIR}, and \texttt{COMSPEC} where present (Windows process creation and parts of the standard library misbehave without them), plus---on Windows only---the two processor-architecture variables discussed in Step 5; and then sets exactly three Python variables: \texttt{PYTHONHASHSEED} to \texttt{0}, \texttt{PYTHONIOENCODING} to \texttt{utf-8}, and \texttt{PYTHONUTF8} to \texttt{1}. Fixing the hash seed removes the one deliberate source of run-to-run nondeterminism CPython injects (randomized string hashing perturbs any code that accidentally depends on set or dict iteration order in unhashed ways); fixing the I/O encoding and UTF-8 mode removes locale-dependent text behavior, so the worker's byte-level output does not vary with the host's regional settings. The execution profile later verifies that the worker environment was \emph{exactly} this constructed mapping---sorted, duplicate-free, subset of the whitelist, value-for-value equal---so an extra inherited variable is not merely unused but structurally rejected (\texttt{E-COMPAT-ENVIRONMENT}).

Why is inherited \texttt{sys.modules} an evidence hazard? Consider the alternative the fresh process replaces: executing the compiler in the calling process, or in a forked copy of it. The calling process is a long-lived research host. By the time a compilation is requested it may have imported hundreds of modules---test frameworks, notebook kernels, plugin systems, earlier versions of CAS code---any of which may have monkey-patched builtins, registered import hooks, installed codecs, or cached state in module globals. \texttt{sys.modules} is a process-global, mutable map, and Python's import system consults it \emph{before} the file system: a module already present under a given name is simply returned, whatever its content. An in-process execution of the hash-verified compiler bytes would therefore run against a module namespace assembled by unhashed history. The receipt would bind the compiler bytes exactly and yet describe almost nothing about what actually executed, because the semantics of every \texttt{import} inside the compiler would be determined by parent-process state that appears nowhere in the evidence. Worse, the contamination is bidirectional and cumulative: each compilation would alter the very state the next compilation runs under, so receipts would not even be comparable with each other. The fresh process resets \texttt{sys.modules} to the interpreter's minimal startup contents, and the receipt records two boolean facts---\texttt{fresh\_process} true, \texttt{parent\_sys\_modules\_inherited} false---that the runtime gate later refuses to pass without (\texttt{DENY\_NONISOLATED\_PYTHON\_PROCESS}).

Two further execution conditions complete Step 4, and each has a rationale of the same kind as the flags.

\emph{The disposable working directory.} The worker's current working directory is a freshly created temporary directory, unique to the launch and removed when the launch completes. The current directory is one of the oldest ambient channels in computing: relative paths resolve against it, some tools consult it for configuration, and---the historical Python-specific hazard---older interpreter configurations put it on the import path, so that a file named like a standard module in the invoker's directory could shadow the real one. The \texttt{-I} flag already removes the import-path exposure; the disposable directory removes the rest. A worker whose current directory is empty, fresh, and about to be deleted can neither find anything there nor leave anything behind that a later run could find: the directory contributes exactly nothing to the execution, which is the only contribution an unhashed location is permitted to make. It also means two compilations never share a working directory, closing the channel by which one run's residue could become another run's input---a small-scale instance of the same reasoning that motivates the fresh process itself.

\emph{The 30-second timeout.} The worker is bounded at 30 seconds of wall clock, enforced by the parent. The bound is not a performance target---the frozen compiler parses a single source unit in milliseconds---but a liveness backstop with an evidence rationale. Without a bound, a worker wedged by any of the host pathologies this chapter and the next catalog (launch storms, file-system stalls, interpreter-initialization anomalies) would hang its compilation, and its caller, indefinitely; the failure would manifest as a silent stall somewhere up the stack rather than as a typed record at the point of fault. With the bound, a wedged worker becomes a worker-start failure with a transport record, attributable and countable like every other failure. The value is generous by three orders of magnitude relative to the expected work precisely so that it never fires on legitimate compilations: a timeout that trips on slow-but-honest runs would create a selection pressure correlated with host load, and Chapter~\ref{ch08:concurrency}'s argument against biased receipt populations applies to timeouts exactly as it applies to retries. Thirty seconds is small enough to keep research pipelines live and large enough that its firing is, in practice, always news about a fault rather than about load.

\begin{ruleenv}[State must be inventoried or absent]\label{ch07:rule:state}
Any interpreter state that can influence the compilation must either appear, content-addressed, in the execution profile, or be structurally absent from the worker. Ambient state that is present but uninventoried is not admitted in any form.
\end{ruleenv}

The worker's module namespace after execution is then \emph{inventoried}: every entry of \texttt{sys.modules} is reported with its origin kind (file, builtin, frozen, or virtual), and every file-backed module must resolve to an absolute path inside the selected Python installation---or be the worker script itself---and must match a recorded digest. A file-backed module outside those roots fails the run with \texttt{E-COMPAT-DEPENDENCY-ESCAPE}. This is the closing of the loop opened by \texttt{-I -S}: the flags make it \emph{unlikely} that anything outside the installation is loaded, and the inventory check makes it \emph{fail closed} if it happens anyway.

\section{Step 5: pre-execution runtime binding without probing}\label{ch07:step5}

The fifth step identifies \emph{what machine and interpreter} the evidence describes---and does so by reading, never by asking.

Before any compiler code executes, the worker samples its runtime identity: the interpreter implementation name (mapped to a display form---\texttt{CPython}, \texttt{PyPy}, and so on---from \texttt{sys.implementation.name}), the three-component version from \texttt{sys.version\_info}, and a platform pair. On Windows the platform name is derived from \texttt{sys.platform} and the machine label from the \texttt{PROCESSOR\_ARCHITEW6432} or \texttt{PROCESSOR\_ARCHITECTURE} environment variables---which is precisely why those two variables are in the environment whitelist of Step 4: they are promoted from ambient state to explicit, profile-hashed \emph{inputs}. On POSIX hosts the pair comes from \texttt{os.uname}. If neither source is available the label is an explicit \texttt{unknown-machine} rather than a guess. Alongside these labels, the worker reports the resolved path, byte count, and raw SHA-256 digest of its own interpreter binary and of the worker script file, and the parent verifies both against fresh reads of the same paths.

What the protocol pointedly does \emph{not} do is probe. It does not query WMI, it does not shell out to \texttt{uname} or \texttt{systeminfo} or \texttt{wmic}, it does not import \texttt{platform} machinery that lazily spawns subprocesses, and it performs no lazy imports at all in the identity path. This restraint is itself a security and evidence decision, for four reasons.

First, \emph{probing executes code that the profile does not bind}. A WMI query traverses a service, a provider stack, and a COM infrastructure; a subprocess probe executes whatever binary the search path resolves. None of those components are content-addressed in the profile, so their outputs would be assertions by unidentified software---exactly the class of input the kernel refuses everywhere else. A runtime label obtained through an unhashed intermediary has the evidentiary standing of hearsay.

Second, \emph{probing widens the attack surface at the worst moment}. The identity sample runs inside the trusted window between process start and compiler execution. Every additional mechanism invoked in that window---service connections, subprocess creation, dynamic library loading---is an additional mechanism whose compromise or misconfiguration can alter the run before the compiler's first instruction, and an additional reason two otherwise identical runs might differ.

Third, \emph{probing is a nondeterminism generator}. WMI calls time out, return transient errors, and can be slow under load; subprocesses inherit scheduling jitter and can fail for reasons unrelated to the machine's identity. The deterministic-closure property that Step 8 will test---two independent executions must agree exactly---would be hostage to infrastructure that fails sporadically. The chosen sources (\texttt{sys} fields, \texttt{os.uname}, two whitelisted environment variables) are plain memory reads within the already-running interpreter: they cannot hang, cannot spawn, and return the same bytes every time on the same host.

Fourth, and most structurally: \emph{the profile must be a function of its recorded inputs}. The execution-profile validator recomputes the expected runtime identity from the same sanitized environment that the worker received and requires the worker's report to match it exactly (\texttt{E-COMPAT-RUNTIME}). That recomputation is only possible because every input to the identity derivation is itself part of the recorded profile. A probe result cannot be validated this way---there is nothing recorded from which to recompute it---so it could only be trusted, and the kernel does not trust.

\begin{proposition}[Runtime identity is replayable]\label{ch07:prop:runtime-replay}
Under the profile's derivation rules, the runtime identity quadruple (implementation, version, platform, machine) is a deterministic function of the interpreter binary and the sanitized environment mapping, both of which are content-addressed in the execution profile. Two runs whose profiles agree on those inputs agree on the quadruple.
\end{proposition}

\begin{proof}[Design argument]
The derivation reads only \texttt{sys.implementation.name}, \texttt{sys.version\_info}, \texttt{sys.platform} or \texttt{os.uname}, and the two whitelisted architecture variables. The first three are constants of the interpreter build identified by the executable digest; the last are values of the recorded environment mapping. No file-system, network, service, or subprocess input participates. The validator enforces agreement by recomputing the quadruple and rejecting mismatches, so a profile asserting a non-derivable identity cannot be constructed.
\end{proof}

The binding is \emph{pre-execution} for the same reason the compiler hash is: an identity sampled after compiler code has run is an identity the compiled code could have influenced. The worker's protocol is strictly staged---verify received compiler digest, sample runtime, execute, inventory, report---and the parent's shape checks on the response (exact key sets at every level, rejected on any deviation) make a worker that reordered or omitted stages structurally distinguishable from a conforming one.

\section{Interlude: the worker wire protocol}\label{ch07:wire}

Steps 4 and 5 described the worker's execution conditions; this interlude describes the conversation between parent and worker, because the wire format is itself part of the evidence discipline and several of its choices block failures that would otherwise be invisible.

The exchange is one request and one response, both single JSON documents. The request travels over standard input and carries exactly five fields: the protocol identifier \texttt{cas\_core\_v8.isolated\_compiler\_worker.v0.1}, the base64-encoded compiler source bytes, the compiler's raw SHA-256, the legacy source text, and the source SHA-256. The request is serialized through the kernel's canonical encoding, so the bytes the worker receives are the same bytes any auditor would reproduce from the recorded inputs. Delivering the compiler over standard input rather than by path is the mechanical realization of Rule~\ref{ch07:rule:no-path-authority}: the worker never opens the compiler file, so no file-system race, symlink trick, or path-resolution surprise can interpose between verification and execution. The worker recomputes the digest of the decoded bytes and refuses on mismatch, giving the pair a two-sided check---the parent verified what it sent, the worker verifies what it received.

The response is read from standard output, with three parsing rules that each close a specific hole.

First, the parent parses the response with a duplicate-key rejecting hook: a JSON object in which the same key appears twice fails the run with \texttt{E-COMPAT-WORKER-PROTOCOL}. Ordinary JSON parsers silently keep the last occurrence, which means two parsers can extract different values from the same bytes---a classic parser-differential channel through which a crafted or corrupted response could tell the shape-checker one thing and the value-consumer another. Rejecting duplicates outright removes the ambiguity instead of adjudicating it.

Second, floating-point literals are parsed directly into exact decimals rather than binary floats. The canonical evidence encoding rejects binary floats; parsing them into \texttt{Decimal} at the wire boundary means a numeric value in a compiler artifact never transits through a lossy binary representation on its way into evidence, and two hosts with different float-formatting behavior cannot produce receipts that differ for that reason alone.

Third, every level of the response is checked against an \emph{exact} key set---not a superset check but set equality. A success response must have exactly the keys \texttt{protocol}, \texttt{ok}, \texttt{artifact}, \texttt{runtime}, \texttt{dependencies}, and \texttt{isolation}; the runtime object exactly its eight keys; the flags object exactly its six; an error response exactly \texttt{protocol}, \texttt{ok}, \texttt{error\_code}, and \texttt{message}. An extra field is as fatal as a missing one. This is unusual---most protocols tolerate unknown fields for forward compatibility---but forward compatibility is precisely what an evidence protocol must not have by accident. An unknown field is unaudited content riding inside a hash-bound record; tolerated today, it becomes load-bearing tomorrow without ever passing review. Version evolution in this kernel happens by changing the protocol identifier, loudly, not by widening what a fixed identifier accepts.

The response's outcome logic is similarly strict in both directions. If the worker reports failure, the parent raises the worker's own typed error code. If the worker reports success but exited nonzero, that contradiction is itself a failure (\texttt{E-COMPAT-WORKER-EXIT}): a process that claims success while signaling failure is not half-trusted on the favorable half. If the worker wrote \emph{anything} to standard error, the run fails with \texttt{E-COMPAT-WORKER-STDERR} even if standard output parsed perfectly---because a conforming worker is silent on stderr, and unexplained diagnostics are unexplained behavior. In every one of these failure paths, the transport receipt of Chapter~\ref{ch08:concurrency}---exit code, byte counts, digests of both streams---is embedded in the error, so the failure record identifies the exact observable behavior without reproducing its content.

\begin{ruleenv}[Exact shapes only]\label{ch07:rule:exact-shapes}
Worker messages are validated by set equality of keys at every level, duplicate keys are rejected, and any deviation---extra field, missing field, wrong protocol identifier, stderr output, success-with-nonzero-exit---fails the compilation closed. There is no lenient parse mode.
\end{ruleenv}

\section{Interlude: freezing the legacy artifact}\label{ch07:freeze}

Between the worker's response and the receipt of Step 7 stands one more transformation that deserves explicit treatment: the compiler's output is \emph{frozen} into exact, immutable evidence values before it is hashed or stored. The freezing pass walks the entire artifact and applies a closed set of rules; anything the rules do not cover fails the run with a typed \texttt{E-COMPAT-ARTIFACT} error.

Null, booleans, and integers pass through unchanged. Strings must satisfy the kernel's canonical-text requirements. Exact decimals must be finite---a NaN or infinity in compiler output is not representable in evidence, so it fails closed rather than serializing as some encoder's best guess. Any binary float that survived to this layer is converted to the exact decimal it denotes, with the same finiteness requirement. Mappings must have string keys, each canonicalized, and are frozen into read-only views; a mapping that uses the encoder's reserved type-tag key is rejected outright, because an artifact that could impersonate the encoding's own tagging machinery could forge type distinctions---the type-pun attack of the encoding chapter, attempted from inside the data. Lists are frozen into immutable list values; bare tuples are rejected with an instruction to use the immutable list type, so that there is exactly one array representation in evidence rather than two that compare unequal. After freezing, the artifact's root must be a mapping whose \texttt{valid} field is a genuine boolean.

Why freeze at all? Two reasons, one about time and one about identity. The time reason: the compilation record outlives the function call that created it, and its digest is computed once. If the artifact remained mutable, any later code touching the record could alter the artifact after digesting, and the record would carry a digest of something it no longer contains---the constructor's re-verification (the recomputed artifact digest must equal the recorded one, else \texttt{E-COMPAT-ARTIFACT-HASH}) would catch this only at reconstruction boundaries, not in a live object. Immutability makes the digest permanently truthful. The identity reason: the freezing rules are exactly the canonical encoding's admission rules, applied eagerly. Freezing at the boundary means a compilation record can never hold a value that the encoder would later reject; there is no artifact that is storable but unhashable, and therefore no evidence object whose identity is undefined.

\begin{remark}
The freezing pass is also where the artifact's diagnostics are derived: the frozen AST is walked once for fractal literals and overloaded RPM calls, producing the typed diagnostic codes of Step 7. Deriving diagnostics from the frozen form, rather than the raw response, guarantees that the diagnostics describe exactly the bytes the digest describes.
\end{remark}

\section{Step 6: advisory-lock serialization of worker launches}\label{ch07:step6}

The sixth step governs \emph{when} workers may start, and it is the protocol's one concession to the messiness of real operating systems.

The empirical problem is concrete: on Windows, when many isolated Python runtimes initialize concurrently, a child process can occasionally exit with a return code and \emph{no} stdout or stderr at all---a transport-level failure with nothing to diagnose. A process-local mutex does not fix this, because independent test runners and research processes overlap on the same host. The protocol therefore serializes worker creation through two nested exclusions: a process-local thread lock, and a cross-process advisory file lock---a single byte locked via \texttt{msvcrt.locking} on Windows and \texttt{fcntl.flock} elsewhere---on a fixed-name file (\texttt{cas-core-v8-compatibility-worker.lock}) in the host temporary directory. The launch policy identifier \texttt{cross-process-exclusive-lock-v1}, the lock path, and all three timing parameters are recorded in every execution profile and verified against the frozen constants at profile construction; a profile asserting different timing is rejected (\texttt{E-COMPAT-ISOLATION}), so the concurrency discipline under which evidence was produced is itself part of the evidence.

The parameters: acquisition is bounded at 90 seconds, after which the compilation fails closed with \texttt{E-COMPAT-WORKER-LOCK}; while waiting, the acquirer polls the non-blocking lock at 10-millisecond intervals; after releasing, the holder sleeps 10 milliseconds before proceeding, as a deliberate scheduler yield so that a process already waiting on the shared lock gets an opportunity to acquire it. Chapter~\ref{ch08:concurrency} develops the rationale for each number and for the advisory (as opposed to mandatory) choice in depth; here we record the two properties that belong to the closure protocol itself.

First, \emph{no retry}. If a worker fails---for any reason, including the very launch-storm failures the lock exists to reduce---the compilation fails. The protocol never launches a replacement worker for a failed one. The full argument is in Chapter~\ref{ch08:concurrency}; the short form is that retrying a nondeterministic failure filters the receipt population by success and thereby converts flakiness into silent selection bias over evidence.

Second, \emph{no fairness claim}. The advisory lock guarantees mutual exclusion of cooperating launchers; it does not guarantee that waiters acquire in order, that any particular waiter is not starved short of its 90-second bound, or anything about non-cooperating processes, which are not obliged to respect an advisory lock at all. The post-release yield is a mitigation, not a guarantee, and the protocol's documentation says so. This honesty is not modesty for its own sake: a fairness property, once asserted, becomes something downstream consumers may build on, and asserting a property the mechanism does not have is exactly the overclaiming failure mode the kernel exists to prevent---applied, this time, to the kernel's own description of itself.

\begin{threat}[Launch-storm evidence corruption]\label{ch07:threat:storm}
Concurrent isolated-interpreter initialization intermittently yields empty-transport child failures. Without serialization, a fraction of compatibility executions fail for reasons uncorrelated with their inputs; without the no-retry rule, that fraction would be silently resampled away, biasing the surviving receipts; without recorded lock parameters, the conditions under which any receipt was produced would be unknowable. The three mechanisms jointly keep the failure visible, bounded, and attributable.
\end{threat}

\section{Step 7: the receipt}\label{ch07:step7}

The seventh step fixes \emph{what is remembered}. Its principle: the receipt must contain every fact a skeptical reviewer needs in order to re-derive the run's identity, and it must state its own limitations as machine-readable fields rather than prose.

The compilation record and the compiler-validation receipt derived from it carry, at minimum:

\emph{Compiler identity, twice.} Both the raw digest of the compiler bytes that executed (which of the two frozen physical variants it was) and the LF-normalized content digest. As argued in Step 1, neither substitutes for the other: the raw digest answers ``which bytes ran,'' the LF digest answers ``which content ran,'' and conflating them has historically hidden both real differences and harmless ones.

\emph{Interpreter and worker bytes.} The resolved path, byte count, and SHA-256 digest of the Python executable, and the same for the worker script. The compiler is not the only code in the room; the interpreter that ran it and the harness that fed it are code too, and a receipt that hashed the compiler while leaving the interpreter anonymous would invite exactly the substitution it closes elsewhere.

\emph{Import search paths.} The worker's complete \texttt{sys.path}, each entry required to be absolute, duplicate-free, and confined to the selected Python installation. The search path is the function from import statements to file contents; recording it makes the import behavior of the run auditable rather than assumed.

\emph{Loaded-module inventory.} Every module present in the worker after execution, with origin kind, path where applicable, and content digest for every file-backed entry; the inventory must be sorted, duplicate-free, and must bind \texttt{\_\_main\_\_} to the worker script's own digest. This is the receipt's answer to ``what code, in total, participated?''---not the code that was supposed to participate, but the code the interpreter actually loaded.

\emph{Exact profiles.} The syntax profile (\texttt{cas-v7.7}), the semantic profile (\texttt{CAS v7.7 operational compatibility only}), the engine profile (\texttt{cas-v7.7/python-only}), and the runtime profile string derived in Step 5. These scope the receipt: any consumer can determine, from the receipt alone, precisely which compatibility claim is being evidenced and which are not.

\emph{Result digest and diagnostics.} The canonical-encoding digest of the frozen compiler artifact, and a typed, nonempty diagnostic tuple. The diagnostics always include the profile marker \texttt{LEGACY\_COMPATIBILITY\_PROFILE}, and add \texttt{UNTYPED\_FRACTAL\_REQUIRES\_BRIDGE} when the artifact contains fractal literals (whose v8 meaning only the explicit bridge of Chapter~\ref{ch08:concurrency} may supply) and \texttt{OVERLOADED\_RPM\_REQUIRES\_TYPE\_CHOICE} when it contains RPM calls whose typing a human must resolve. Diagnostics are how the receipt marks the places where legacy content touches questions the profile refuses to answer.

\emph{The research boundary.} Every record embeds the Ring-4 boundary triple---admission \textsc{open}, product \textsc{not licensed}, disclosure false---and every constructor in the chain re-verifies it. A receipt whose boundary is anything else cannot be built.

\emph{Capability booleans, including the false ones.} The receipt asserts \texttt{isolated\_mode}, \texttt{ignore\_environment}, \texttt{fresh\_process} true and \texttt{parent\_sys\_modules\_inherited} false, as established in Step 4---and it asserts \texttt{os\_sandboxed} false and \texttt{host\_os\_dependencies\_content\_addressed} false, as Section~\ref{ch07:not-sandbox} discusses. It also carries an explicit self-hash scope string, \texttt{CORRUPTION\_DETECTION\_ONLY\_NOT\_AUTHENTICATION}, so that the receipt's own integrity digest cannot be quietly promoted into an authentication claim even by documentation drift.

\begin{ruleenv}[Receipts state their own weaknesses]\label{ch07:rule:weakness-fields}
A capability the execution does not have is recorded as an explicit false field, verified at construction (a receipt asserting it true is rejected), and re-checked at assessment (a receipt asserting it true is denied with \texttt{DENY\_FALSE\_OS\_ISOLATION\_CLAIM}). Limitations are data, not caveats.
\end{ruleenv}

The design debt here is to provenance practice \cite{prov}: the receipt separates the entity (source and compiler bytes), the activity (the worker execution, with its policy and timing), and the environment of the activity (interpreter, modules, paths), rather than collapsing them into a single opaque success bit. And the receipt's whole content is serialized through the canonical encoding of the earlier chapters \cite{jcs}, so its digest is stable and its comparison semantics are exact.

\section{Step 8: fresh replay before any structural decision}\label{ch07:step8}

The eighth step determines \emph{when the evidence may be used}: never on faith in a stored record, always after fresh re-derivation.

The requirement operates at two nested levels. At the compilation level, every compile is already a determinism experiment: the protocol runs \emph{two} independent fresh workers over the same retained source and compiler bytes and requires their frozen artifacts to have equal canonical digests (\texttt{E-COMPAT-NONDETERMINISTIC-RESULT} otherwise) and their execution profiles to have equal digests (\texttt{E-COMPAT-NONDETERMINISTIC-PROFILE}, with the bounded diagnostic of Chapter~\ref{ch08:concurrency}). A compilation record thus never represents a single run's say-so; it represents an agreement.

At the decision level, the requirement bites harder. When a compilation record is later presented as evidence for a structural decision---in the assessment and gate machinery of Chapter~\ref{ch09:decision}---the kernel does not compare the stored record against stored expectations. It first re-verifies that every file the record's execution profile binds (interpreter, worker, file-backed modules) still has its recorded bytes (\texttt{E-COMPAT-DEPENDENCY-RACE} on any drift), and then \emph{re-executes the whole protocol}: same retained source bytes, same retained compiler bytes, two fresh workers, full verification, and requires the replay's complete canonical serialization---artifact, profile, diagnostics, everything---to be byte-identical to the supplied record (\texttt{E-COMPAT-REPLAY-MISMATCH} otherwise). Only the replay's receipt, not the stored one, feeds the decision.

What does this block? The stored-receipt attack surface in its entirety. A receipt is bytes; bytes can be fabricated, replayed from another host, generated by a tampered earlier run, or produced honestly under conditions that no longer hold. Signature schemes address some of this at the cost of a key-management trust root the kernel deliberately does not yet have. Fresh replay addresses it differently: the receipt is demoted from \emph{proof} to \emph{prediction}. Presenting a compilation record to the gate amounts to predicting, in exact canonical bytes, what a fresh deterministic-closure execution will produce; the gate runs the experiment and checks the prediction. A fabricated receipt fails because the fabricator must predict digests of an execution they cannot run in advance on the verifier's host; a stale receipt fails because some bound file has drifted; an honest receipt passes and, in passing, demonstrates that the claimed determinism actually holds \emph{now}, on \emph{this} machine, which is the only fact the decision needs.

\begin{invariant}[No decision on stored evidence]\label{ch07:inv:replay}
No structural decision in the kernel is reachable from a stored compatibility record alone. Every path to a structural pass includes a fresh execution of the full closure protocol whose canonical serialization must equal the record byte-for-byte.
\end{invariant}

The cost is real: replay doubles-and-more the execution work, and confines structural decisions to hosts that hold the same interpreter and compiler bytes. The kernel accepts this cost knowingly. In a research program whose documented failure history includes self-generated artifacts masquerading as authority, cheap trust in stored records is precisely the resource an overclaim consumes. Making evidence expensive to \emph{assert} and cheap to \emph{check} is the correct asymmetry; replay makes assertion and checking the same act.

\section{The isolation boundary: what this protocol is not}\label{ch07:not-sandbox}

Everything above concerns Python-process isolation, and it is important to state exactly what that phrase does and does not mean. The worker is a fresh operating-system process with a sanitized environment, an isolated interpreter configuration, an inventoried module namespace, and a bounded lifetime. It is \emph{not} sandboxed. It runs as the invoking user, with the invoking user's privileges, on the shared file system, with no syscall filter, no capability restriction, no network namespace, no resource jail beyond the 30-second timeout. A hostile compiler---which the frozen, hash-bound compiler is not, but the point of an evidence kernel is not to need that assumption---could read files, write files, or open sockets, and nothing in this protocol would prevent it. The protocol constrains and records the \emph{Python} face of the execution; the operating-system face is constrained only by the host's ordinary rules.

Likewise, the content-addressing perimeter has a definite edge. Inside it: the source, the compiler, the interpreter binary, the worker script, every file-backed Python module, the environment mapping, the search paths, the launch policy. Outside it: the native libraries the interpreter loads (the C runtime, OpenSSL, the OS loader's dependency graph), the kernel and system services beneath them, and the hardware. Two receipts with identical profiles can therefore have been produced above different native substrates, and the receipt cannot tell you. For the deterministic-closure property this residual matters less than it might---the replay requirement re-runs on the actual current substrate, so a substrate-dependent divergence surfaces as a replay mismatch rather than passing silently---but for any claim of the form ``this execution was hermetic,'' the receipt is not evidence, and it says so itself.

\begin{boundary}[Non-sandbox, machine-readable]\label{ch07:bd:nonsandbox}
The profile's \texttt{os\_sandboxed} and \texttt{host\_os\_dependencies\_content\_addressed} fields are constant false in v0.1. The constructors reject a profile or receipt asserting either true (\texttt{E-COMPAT-ISOLATION}); the assessment denies such a receipt (\texttt{DENY\_FALSE\_OS\_ISOLATION\_CLAIM}); and every structural pass carries the two limitation codes \texttt{LIMITATION\_NO\_OS\_SANDBOX} and \texttt{LIMITATION\_HOST\_OS\_DEPENDENCIES\_NOT\_CONTENT\_ADDRESSED} in its decision record. A downstream consumer cannot honestly read any artifact of this protocol as hermetic-execution evidence, because the artifact itself denies it in typed fields.
\end{boundary}

This is an unusual design choice and worth dwelling on. Most systems describe their guarantees; few encode their non-guarantees as validated, constant, machine-checked fields. The kernel does so because the alternative---limitations living in documentation---has a known failure mode: documentation and behavior drift apart, and the strongest claim anyone remembers becomes the claim in circulation. A false capability field cannot drift. If a future version of the kernel gains a real OS sandbox, admitting it will require changing validators that currently \emph{reject} the true value---a loud, reviewable act---rather than quietly starting to say yes. The threat being modeled here is not an external attacker but the program's own future enthusiasm \cite{shostack}: the boundary field is a precommitment device against the kernel's maintainers, in the same spirit as the constant-false authorization fields of Chapter~\ref{ch09:decision}.

\section{Reading a receipt: a reviewer's walkthrough}\label{ch07:walkthrough}

The steps above are stated mechanism-first. It is worth walking the same ground once more from the opposite direction: a reviewer holds a compiler-validation receipt and wants to know what, exactly, they may conclude from it. The exercise makes the protocol's division of labor concrete and shows where each step's guarantee is consumed.

The reviewer's first question is \emph{scope}: what kind of claim is this receipt even attempting? The answer is printed in the receipt itself. The receipt version string identifies the record contract; the syntax, semantic, and engine profile fields identify the claim family as CAS v7.7 operational compatibility under the Python-only engine; the \texttt{legacy\_only} flag (validated true) and the structural absence of any v8 translation field establish that no v8 meaning is being asserted; and the embedded Ring-4 boundary establishes that whatever the receipt shows, it is research evidence---admission open, product not licensed, disclosure withheld. A receipt cannot soft-pedal any of this, because each field is validated at construction and re-checked at assessment; the reviewer does not need to trust the surrounding prose, which is the point.

The second question is \emph{identity}: what ran? The receipt names the source digest, both compiler digests, the interpreter digest, and the worker digest. The reviewer can check the compiler digests against the frozen constants of Definition~\ref{ch07:def:compiler-identity} independently---the constants are published in the Gate~4 record---and can distinguish, from the raw digest, whether the CRLF or LF physical variant executed. If the reviewer holds the same artifacts, every one of these digests is recomputable from bytes; nothing requires taking the receipt's word.

The third question is \emph{conditions}: under what regime did it run? The receipt's boolean block answers with validated facts rather than assurances: isolated mode, environment ignored, fresh process, no inherited module state---and, in the false fields, no OS sandbox and no host-OS content addressing. The execution profile behind the receipt adds the sanitized environment (names and values from the fixed whitelist), the confined search paths, the full module inventory, and the launch policy with its exact timing constants. A reviewer who suspects that some receipt was produced under looser conditions has a decidable question, not a judgment call: looser conditions cannot satisfy the profile validators, so either the receipt's digests do not verify, or the conditions were as stated.

The fourth question is \emph{meaning}: what did the run produce, and what interpretive cautions attach? The result digest identifies the frozen artifact exactly; the diagnostics enumerate the caution flags. If \texttt{UNTYPED\_FRACTAL\_REQUIRES\_BRIDGE} is present, the artifact contains fractal digit sequences whose v8 interpretation the receipt explicitly does not supply---only the typed bridge of Chapter~\ref{ch08:concurrency}, with a caller-owned catalog, may do that, and even then only as type-correct composition, never digit meaning. If \texttt{OVERLOADED\_RPM\_REQUIRES\_TYPE\_CHOICE} is present, a human type decision is pending and no downstream automation should pretend otherwise.

The fifth and decisive question is \emph{currency}: is any of this still true, here, now? Here the reviewer's answer is procedural rather than inspective: present the record to the assessment machinery and let Step 8 replay it. If the replay passes, the reviewer has learned something no stored receipt could tell them---that the deterministic closure holds on the present host at the present moment. If it fails, the failure code says why: drifted dependency bytes, a changed interpreter, a nondeterministic compiler behavior surfacing, or a record that never corresponded to a real execution at all. In both outcomes the reviewer's conclusion rests on a fresh experiment, and the stored receipt served only as the experiment's exact specification.

\begin{remark}
Note what is absent from the walkthrough: at no point does the reviewer ask \emph{who} produced the receipt. This is deliberate and double-edged. The protocol makes receipts verifiable without provenance of authorship---anyone with the bytes can check everything---but by the same token a verified receipt carries no authorship claim, no review status, and no authority. The question ``who vouches for this?'' has no answer inside the compatibility profile, and Chapter~\ref{ch09:decision} shows that the kernel answers it, for now, with a constant: no one.
\end{remark}

\section{The protocol as a whole}\label{ch07:whole}

Read together, the eight steps have a shape worth making explicit. Steps 1--3 fix the \emph{inputs}: what runs (hash-verified source and compiler bytes), under what regime (one frozen engine profile), and how much (one source unit, no imports). Steps 4--6 fix the \emph{execution}: where (a fresh, flag-isolated, environment-sanitized process), on what identified substrate (pre-execution runtime binding without probing), and when (serialized launches, bounded waits, no retries). Steps 7--8 fix the \emph{output}: what is remembered (a receipt that re-derives, scopes, and self-limits) and when it may be believed (only upon fresh replay).

\begin{finding}[Closure, in the profile's honest sense]\label{ch07:finding:closure}
Within the recorded perimeter---source, compiler, interpreter, worker, Python modules, environment, search paths, launch discipline---a compatibility execution is a deterministic function, and the protocol tests this determinism on every compile (two workers) and again at every decision (fresh replay). Outside the perimeter---native libraries, operating system, hardware---no closure is claimed, and the receipts carry that non-claim as data. Both the Gate~4 test evidence for this behavior and the profile itself are Ring-4 research artifacts: structural software receipts under exact fixtures, not certification, review, or authority.
\end{finding}

The next chapter takes up the two components of this protocol that most repay closer analysis---the concurrency discipline of Step 6 and the failure taxonomy that Steps 4--6 feed---together with the one bridge by which legacy artifacts may touch v8 structure at all.
